!TeX root = ManuscriptGuide.tex
\newpage
\clearpage
\section{Manuscript Organization Guide}
\begin{verbatim}
=========================================================
Manuscript Organization Guide
=========================================================

Purpose of this Guide

This guide provides a structural overview of the manuscript.
It identifies the purpose of each section, the questions
addressed, the concepts introduced, the figures and tables
used to support those concepts, and the relationships among
sections.

The guide is intended to improve transparency of the paper's
organizational structure and facilitate navigation by both
human readers and AI-assisted document-analysis tools.

=========================================================
Introduction
=========================================================

Purpose

    Introduce the collision-detection problem,
    identify the remaining architectural
    limitation in existing GPU approaches,
    and present the gPCD framework.

Questions Answered

    What are particle methods?

    Why do particle counts become large?

    Why is collision detection difficult?

    How have GPUs been used to address this
    challenge?

    What architectural limitation remains?

    What is gPCD?

    What are the contributions of this work?

Key Concepts Introduced

    - Particle methods
    - Large-scale particle simulation
    - Broad-phase collision detection
    - GPU acceleration
    - Architectural limitation
    - Graphics-pipeline occupancy generation

Figures/Tables

    None

Inputs

    None

Outputs

    - Architectural limitation
    - gPCD framework
    - Contributions of the paper

Key Takeaway

    Existing GPU collision-detection methods
    accelerate processing, but occupancy
    generation remains largely independent of
    rendering. The gPCD framework addresses
    this limitation by integrating occupancy
    generation directly into the graphics
    pipeline.

=========================================================
Background and Related Work
=========================================================

Purpose

    Review representative collision-detection
    methodologies, execution architectures,
    and rendering relationships that motivate
    the gPCD framework.

Questions Answered

    What collision-detection methodologies
    exist?

    How have collision-detection systems been
    implemented?

    What is the relationship between rendering
    and collision detection?

    What architectural gap remains?

Key Concepts Introduced

    - Broad-phase methodologies
    - Uniform grids
    - Spatial hashing
    - Hierarchical subdivision
    - Sweep and prune
    - Image-space methods
    - CPU-resident architectures
    - GPU-assisted architectures
    - GPU-resident architectures
    - Rendering as a consumer of simulation
      state
    - Architectural gap

Figures/Tables

    None

Inputs

    - Architectural limitation identified in
      the Introduction

Outputs

    - Architectural gap definition
    - Motivation for graphics-pipeline
      occupancy generation

Key Takeaway

    Existing methodologies and architectures
    accelerate collision detection, but
    occupancy generation remains largely
    independent of the rendering pipeline.

=========================================================
Methodology
=========================================================

Purpose

    Describe the architecture and operation of
    the gPCD framework.

Questions Answered

    What is the overall gPCD architecture?

    How is occupancy generated?

    How are cells identified?

    How are occupancy lists constructed?

    What is the potentially colliding set?

    How are collisions evaluated?

    How are duplicate collisions prevented?

Key Concepts Introduced

    - Graphics-pipeline occupancy generation
    - Broad phase
    - Narrow phase
    - AABB corner evaluation
    - Cell addressing
    - Flattened cell indices
    - Particle corner occupancy list
    - Cell-array particle occupancy list
    - Duplicate suppression
    - Potentially colliding set (PCS)
    - Source-owned collision evaluation

Figures/Tables

    Figure:
        overview_architecture

    Figure:
        CellSpans

Inputs

    - Architectural gap from Background

Outputs

    - Occupancy-generation workflow
    - PCS construction
    - Source-owned collision evaluation
    - Data structures used throughout the
      framework

Key Takeaway

    Broad-phase occupancy generation is
    integrated directly into the graphics
    pipeline, producing the PCS used by the
    narrow phase while rendering and occupancy
    construction proceed concurrently.

=========================================================
Complexity Analysis
=========================================================

Purpose

    Derive the computational and storage
    complexity of the gPCD framework.

Questions Answered

    What is the complexity of occupancy
    generation?

    What is the complexity of narrow-phase
    evaluation?

    What is the worst-case interaction count?

    What are the storage requirements?

Key Concepts Introduced

    - O(N) occupancy generation
    - O(K) narrow-phase evaluation
    - O(N+K) total complexity
    - O(N) particle storage
    - O(S^3W) cell storage

Figures/Tables

    Table:
        Complexity Comparison

Inputs

    - Occupancy-generation workflow
    - PCS construction
    - Narrow-phase evaluation

Outputs

    - Complexity bounds
    - Storage bounds
    - Theoretical expectations for scaling

Key Takeaway

    The framework exhibits O(N) occupancy
    generation and O(N+K) overall complexity,
    where K represents candidate interactions
    contained within the PCS.

=========================================================
Experimental Setup
=========================================================

Purpose

    Define the hardware platform, benchmark
    configurations, verification procedures,
    and performance metrics used to evaluate
    the framework.

Questions Answered

    What hardware was used?

    How was the framework implemented?

    What limits the maximum particle count?

    How were occupancy lists sized?

    What benchmark configurations were used?

    How was correctness verified?

    What metrics were measured?

Key Concepts Introduced

    - Hardware platform
    - Vulkan implementation
    - Storage-buffer constraints
    - Occupancy-list sizing
    - Verification mode
    - Performance mode
    - Throughput
    - Time-per-particle

Figures/Tables

    None

Inputs

    - Complexity expectations
    - Framework implementation

Outputs

    - Benchmark methodology
    - Measurement definitions
    - Verification methodology

Key Takeaway

    Performance measurements are obtained from
    verified collision-detection results using
    benchmark configurations that remain within
    the storage and occupancy constraints of
    the evaluation platform.

=========================================================
Experimental Results
=========================================================

Purpose

    Present measured performance of the gPCD
    framework.

Questions Answered

    How does memory locality affect
    performance?

    What throughput is achieved?

    How does time-per-particle behave?

    How does the framework scale?

    How large are deviations from ideal
    linear behavior?

Key Concepts Introduced

    - Memory-locality benchmark
    - Throughput
    - Time-per-particle
    - Scaling exponents
    - Curvature bounds
    - Throughput-transition regime
    - Stable large-scale regime

Figures/Tables

    Memory Locality Benchmark

    Performance Summary

    Throughput and Time-per-Particle

    Log-Log Regression

    Curvature Analysis

Inputs

    - Experimental setup
    - Performance metrics

Outputs

    - Throughput measurements
    - Scaling measurements
    - Locality measurements
    - Curvature measurements

Key Takeaway

    The framework demonstrates high throughput
    and near-linear scaling across the
    evaluated particle range while supporting
    multi-million-particle collision detection
    on a single CPU-GPU system.

=========================================================
Discussion
=========================================================

Purpose

    Interpret the architectural significance,
    performance characteristics, and practical
    implications of the framework.

Questions Answered

    Why is gPCD architecturally different?

    What synchronization advantages result
    from the architecture?

    Do the results support the complexity
    analysis?

    What role does memory locality play?

    What practical particle-system scales are
    achievable?

    What visualization advantages result from
    the framework?

Key Concepts Introduced

    - Architectural significance
    - Source-owned synchronization
    - Scaling interpretation
    - Memory locality effects
    - Practical particle-system scale
    - GPU-resident visualization

Figures/Tables

    Figure:
        DensityGraphs

Inputs

    - Experimental results

Outputs

    - Interpretation of results
    - Practical implications
    - Visualization implications

Key Takeaway

    The measured performance is consistent with
    the complexity analysis and demonstrates
    that graphics-pipeline occupancy generation
    is a viable architecture for large-scale
    particle collision detection.

=========================================================
Conclusion and Future Work
=========================================================

Purpose

    Summarize the contributions, findings, and
    future research directions.

Questions Answered

    What was developed?

    What was demonstrated?

    Why is it significant?

    What future work is planned?

Key Concepts Summarized

    - Graphics-pipeline occupancy generation
    - Near-linear scaling
    - Multi-million-particle collision
      detection
    - GPU-resident architecture

Figures/Tables

    None

Inputs

    - Discussion

Outputs

    - Final conclusions
    - Future research directions

Key Takeaway

    The gPCD framework demonstrates that
    graphics-pipeline occupancy generation can
    support large-scale particle collision
    detection while maintaining high throughput
    and near-linear scaling on a single
    CPU-GPU system.

=========================================================
End of Manuscript Organization Guide
=========================================================
\end{verbatim}